\documentclass[11pt]{scrartcl}

\title{Projektplan: SDDC}
\author{Silvan Adrian \\ Fabian Binna}
\date{\today{}}

\usepackage[ngerman]{babel}
\usepackage[automark]{scrpage2}
\usepackage[colorlinks = true,
linkcolor = black]{hyperref}
\usepackage{color}
\usepackage[normalem]{ulem}
\usepackage{scrpage2}
\usepackage{graphicx}
\usepackage{tabularx}
\graphicspath{ {./images/} }
\usepackage{tikz}
\usetikzlibrary{positioning,shadings}
\usetikzlibrary{arrows}
\pagestyle{scrheadings}

\definecolor{logocolor}{RGB}{6,103,236}

\clearscrheadfoot
\ihead{\includegraphics[scale=0.3]{SDDC}}
\ohead{Projekt: SDDC}
\ifoot{Projektplan: SDDC}
\cfoot{Version: 1.01}
\ofoot{Datum: \today{}}
\setheadsepline{0.5pt}
\setfootsepline{0.5pt}

\usepackage{ucs}
\usepackage[utf8]{inputenc}
\usepackage[T1]{fontenc}

\begin{document}
\def\arraystretch{1.5}
\begin{titlepage}
\begin{center}
\vspace{10em}
\includegraphics[scale=2]{SDDC}
\vspace{10em}
\end{center}
\begin{center}
\huge {Projektplan}
\end{center}
\begin{center}
\vspace{10em}
\LARGE {Silvan Adrian} \\
\LARGE {Fabian Binna}
\end{center}

\end{titlepage}

\newpage
\section{Änderungshistorie}
\begin{tabularx}{\textwidth}{l l X l}
\textbf{Datum} & \textbf{Version} & \textbf{Änderung}  & \textbf{Autor} \\
\hline
\textbf{17.09.15} & 1.00 & Erstellung des Dokuments & Gruppe \\
\textbf{18.09.15} & 1.01 & Organisationsstruktur + Einführung + 
Projektübersicht & Silvan Adrian\\

\end{tabularx}

\newpage
\tableofcontents
\newpage

\section{Einführung}
\subsection{Zweck}
Dieses Dokument beschreibt die Planung der SA SDDC.
\subsection{Gültigkeitsbereich}
Dieses Dokument ist während des ganzen Projekts gültig und wird laufend aktualisiert.

\subsection{Referenzen}
-
\section{Projektübersicht}
Es soll eine generische API oder bestehende API verwendet werden, um über ein 
Dashboard direkt Services zu abonnieren, dabei wird für jede Virtuelle Maschine 
Compute, Storage und Network benötigt.

\subsection{Zweck und Ziel}
 Gelerntes anwenden und sich mit der gestellten Aufgabenstellung auseinandersetzen 
 und dabei neues lernen.


\subsection{Primäre Features}

\subsection{Erweiterte Features}

\subsection{Lieferumfang}


\subsection{Annahmen und Einschränkungen}

\section{Projektorganisation}
Wir setzen beim Project auf Scrum, wodurch lediglich die Projektrollen Product 
Owner, Entwickler und Scrum-Master vorhanden sind.
Es gibt keine Unterscheidung zwischen speziellen Fachgebieten und das Team 
als ganzes ist gefragt.
Daily Meetings (ca. 10 - 15min) werden während der Pause oder über Skype stattfinden, um sich 
gegenseitig über Probleme, abgeschlossene Arbeitspakete etc. zu informieren.
Nach jedem Abschluss eines Sprints wird ein Sprintreview geplant, um zu prüfen 
ob alles gemäss Plan erledigt werden konnte + die neuste Version der Software 
released.
\subsection{Organisationsstruktur}
\begin{tikzpicture}[auto, node distance = 0.4cm, thick,
  every node/.style = {rectangle, font = \sffamily, white,
    top color = logocolor!90!black, bottom color = logocolor!60!black,
    text width = 3cm, align = center, minimum height = 1cm}]
  \node (PO) {\textbf{Beat Stettler} \\
      \vspace{0.2cm}
  \includegraphics[width=0.8\textwidth]{stettler}
   \vspace{0.2cm}
  Product Owner};
  \coordinate [below = 3cm of PO] (Mitte);
  \node (Silvan) [left  = 2mm of Mitte]  {\textbf{Silvan Adrian} \\
    \vspace{0.2cm}
  \includegraphics[width=0.8\textwidth]{silvan}
   \vspace{0.2cm}
Entwickler};
  \node (Fabian) [right = 2mm of Mitte]  {\textbf{Fabian Binna} \\
  \vspace{0.2cm}
  \includegraphics[width=0.8\textwidth]{fabian}
   \vspace{0.2cm}
 Entwickler};
  \draw [logocolor!60!black,thick]
    (PO)    -- (Mitte) -- (Silvan)
    (Fabian)   -- (Mitte); 
\end{tikzpicture}

\subsection{Externe Schnittstellen}
\textbf{Projektbetreuer:} \\
Beat Stettler (\href{mailto:Beat.Stettler@ins.hsr.ch}{Beat.Stettler@ins.hsr.ch})  
\\
Urs Baumann (\href{mailto:Urs.Baumann@ins.hsr.ch}{Urs.Baumann@ins.hsr.ch})


\section{Managment Abläufe}

\subsection{Kostenvoranschlag}


\subsection{Sprints}

\subsection{Besprechungen}



\section{Risikomanagement}

\subsection{Risiken}



\subsection{Umgang mit Risiken}

\section{Arbeitspakete}

Die Arbeitspakete werden in Open Project erstellt und gepflegt.
Lesender Zugriff ist anonym möglich, schreibender nur eingeloggt (Projekt ist öffentlich). 
\\
Link zur Open Project Instanz : \href{http://sddc.silvn.com}{http://sddc.silvn.com}

\section{Infrastruktur}

\subsection{Entwicklungsinfrastruktur}

\begin{tabularx}{\textwidth}{c c c c}
\textbf{Name} & \textbf{Hardware} & \textbf{Betriebssystem} & \textbf{IDE} \\
\hline
Silvan Adrian & MacBook Pro & OSX 10.10.5 &  N/A \\
\hline
Fabian Binna & Lenovo T530s & Windows 10 &  N/A\\
\hline
\end{tabularx}

\subsection{Tools/Software}
\begin{itemize}
    \item \textbf{BuildServer:} Travis-CI
    \item \textbf{Konfigurationsmanagement:} GIT
    \item \textbf{Notifications:} Slack
\end{itemize}

\subsection{Kommunikationsmittel}
\begin{itemize}
    \item E-Mail
    \item Skype
    \item Open Project (Kommentare)
    \item GitHub (Issues,Kommentare)
    \item Whatsapp
    \item Slack
\end{itemize}

\section{Qualitätsmassnahmen}

\subsection{Dokumentation}

\subsection{Projektmanagement}

\subsection{Entwicklung}

\subsubsection{Vorgehen}

\subsubsection{Unit Testing}

\subsubsection{Code Review}

\subsubsection{Code Style Guidelines}

\subsection{Testen}

\subsubsection{Komponententest}

\subsection{Systemtest}

\end{document}